\reading{LTR-12}{Managing Non-Deposit Liabilities}{STRIKE FIRST}{2}

\begin{lrkeybox}[Learning objectives — official 2026 spine wording]
\begin{enumerate}[label=\textbf{\alph*.}]
\item Distinguish between the various sources of non-deposit liabilities at a bank.
\item Describe and calculate the available funds gap.
\item Discuss factors affecting the choice of non-deposit funding sources.
\item Calculate overall cost of funds using both the historical average cost approach and the pooled-funds approach.
\end{enumerate}
{\footnotesize\color{FRMGrey}Source: Peter Rose and Sylvia Hudgins, \textit{Bank Management \& Financial Services}, 9th Edition, Chapter 13.}
\end{lrkeybox}

\begin{lrtrapbox}[Both calculation objectives are missing from the source notes]
Objective \textbf{b} ends with the words ``It is calculated as'' and then stops --- there is no formula on the page and no image carrying one. Objective \textbf{d} has a heading and \textbf{no content at all}. Those are the two computational objectives in a \textbf{STRIKE FIRST} reading, and one of them (AFG) sits on the \S5B.10 formula-recall watchlist. Both are gap-filled below from the GARP curriculum.

The Crash layer also labels this reading \texttt{76.a--d}; the Lecture layer uses \texttt{76} for \textbf{LTR-17}. Renumbered to the spine.
\end{lrtrapbox}

% =====================================================================
\lo{LTR-12 a}{Distinguish between the various sources of non-deposit liabilities at a bank.}

\begin{center}
\scriptsize
\begin{longtable}{@{}p{3.5cm} p{10.6cm}@{}}
\toprule
\rowcolor{FRMSoft}\textbf{Source} & \textbf{What it is} \\
\midrule
\endfirsthead
\toprule
\rowcolor{FRMSoft}\textbf{Source} & \textbf{What it is} \\
\midrule
\endhead
\textbf{1. Federal funds} &
Short-term borrowings between banks, to meet reserve requirements, clear cheques, or buy Treasury securities. \\
\addlinespace
\textbf{2. Repurchase agreements} &
Short-term \textbf{secured} borrowings. The borrower sells high-quality securities to get cash, then repurchases them at maturity with interest. \\
\addlinespace
\textbf{3. Discount window borrowing} &
Short-term loans from the Federal Reserve to banks with acceptable collateral, such as Treasury securities. Three types:
\begin{itemize}
\item \textbf{Primary credit} --- overnight loans at a \textbf{low} rate.
\item \textbf{Secondary credit} --- for institutions not qualifying for primary credit, at a \textbf{higher} rate.
\item \textbf{Seasonal credit} --- longer-term loans for seasonal needs, at an \textbf{intermediate} rate.
\end{itemize} \\
\addlinespace
\textbf{4. Federal Home Loan Bank} &
Provides loans to mortgage lenders, with home mortgages as collateral. Offers \textbf{below-market} interest rates. \\
\addlinespace
\textbf{5. Negotiable CDs} &
Large (\$100,000+) CDs, functioning similarly to fed funds and repos for short-term funding. Four main types:
\begin{itemize}
\item \textbf{Domestic CDs} --- issued by US institutions inside the territory of the United States.
\item \textbf{EuroCDs} --- dollar-denominated CDs issued by banks \emph{outside} the United States.
\item \textbf{Yankee CDs} --- sold by the largest foreign banks active in the United States (Deutsche Bank, HSBC) through their \emph{US branches}.
\item \textbf{Thrift CDs} --- sold by nonbank savings institutions.
\end{itemize} \\
\addlinespace
\textbf{6. Eurocurrency deposits} &
Dollar-denominated deposits in banks outside the United States. \\
\addlinespace
\textbf{7. Commercial paper} &
Short-term notes, maturities normally from three or four days to nine months, issued by well-known companies to raise working capital. Generally sold \textbf{at a discount from face value}, through security dealers or by direct contact with the issuing company. \\
\addlinespace
\textbf{8. Long-term non-deposit funds} &
Borrowing in the long-term debt market. Usually viewed as a source of \textbf{secondary capital} rather than a source of funding. \\
\bottomrule
\end{longtable}
\end{center}

\begin{lrtrapbox}[Two sibling sets built for swapping]
\begin{itemize}
\item \textbf{EuroCD versus Yankee CD.} A EuroCD is a dollar CD issued \emph{outside} the US; a Yankee CD is issued by a \emph{foreign} bank \emph{inside} the US. Both involve a foreign element and the direction is the only discriminator.
\item \textbf{The three discount window rates.} Primary is \textbf{lowest}, secondary is \textbf{highest}, seasonal is \textbf{intermediate}. Note that this is an internal ordering; relative to the wider market the \textbf{primary credit rate is generally among the highest short-term borrowing rates}, because Federal Reserve credit is priced \emph{above} the central bank's target for the fed funds rate. Both facts are true and they look contradictory, which is exactly why they make good distractor material.
\item \textbf{Long-term debt is secondary capital, not funding.} An option treating it as a funding source of first resort has missed the reading's framing.
\end{itemize}
\end{lrtrapbox}

% =====================================================================
\lo{LTR-12 b}{Describe and calculate the available funds gap.}

The demand for non-deposit funds is determined basically by the size of the gap between the institution's total credit demands and its deposits and other available monies. Management must be prepared to meet not only today's credit requests but all those it can reasonably anticipate --- so projections must be based on knowledge of the current and probable future funding needs of each institution's customers, especially its largest borrowers.

\gapbadge
\gapnote{The source notes define the gap in words and then stop at ``It is calculated as''. The formula, its worked example and the treatment below are written from the GARP curriculum (Rose and Hudgins, Chapter 13, equation 13.3).}

\begin{lrgapbox}
The objective says \emph{describe and calculate}. The definition without the equation covers half of it, and the half that is missing is the one on the \S5B.10 formula-recall watchlist.
\end{lrgapbox}

\begin{lrfmlbox}[Available funds gap]
\begin{align*}
\text{AFG} \;=\;& \text{Current and projected loans and investments the institution desires to make} \\
&-\; \text{Current and expected deposit inflows and other available funds}
\end{align*}
\vk{A positive AFG is a \emph{funding shortfall} to be covered from non-deposit sources. The asset side of the subtraction includes not only new loans but planned securities purchases \textbf{and} expected drawings on existing credit lines.}
\end{lrfmlbox}

\begin{lrexbox}[Estimating next week's available funds gap]
A commercial bank has new loan requests meeting its quality standards of \$150 million. It wishes to purchase \$75 million in new Treasury securities being issued this week, and it expects drawings on credit lines from its best corporate customers of \$135 million. Deposits and other customer funds received today total \$185 million, and those expected in the coming week will bring in another \$100 million. Estimate the available funds gap for the coming week.
\tcblower
\textbf{Step 1 --- total the uses of funds.}
\[ 150 + 75 + 135 = \$360 \text{ million} \]
All three are things the bank intends to fund: new loans, new securities, and \emph{drawings on lines already committed}.

\textbf{Step 2 --- total the sources of funds.}
\[ 185 + 100 = \$285 \text{ million} \]

\textbf{Step 3 --- take the difference.}
\[ \text{AFG} = 360 - 285 = \mathbf{\$75 \text{ million}} \]

\textbf{What this means, and what happens next.} \$75 million must be raised from non-deposit sources. In practice \textbf{most institutions will add a small amount to this estimate}, to cover unexpected credit demands or unanticipated shortfalls in deposits and other inflowing funds. The AFG is a planning figure, not a floor.
\end{lrexbox}

\begin{lrtrapbox}[Three ways the AFG is corrupted]
\begin{itemize}
\item \textbf{Dropping the credit-line drawings.} $150 + 75 - 285 = -\$60$ million, a clean sign-flipped answer that reads as a \emph{surplus}. Expected drawings on committed lines are a use of funds even though no new lending decision is being made.
\item \textbf{Dropping the securities purchase.} $150 + 135 - 285 = \$0$ --- and \S5B.4 reminds us that \textbf{zero is an offered answer} in GARP questions.
\item \textbf{Counting only funds already received.} $360 - 185 = \$175$ million. Both current \emph{and} expected inflows belong in the subtraction.
\end{itemize}
Note also that the GARP text's own worked line prints the sum as ``\$306'' where the addition gives \$360; the stated answer of \$75 confirms \$360 is intended. A transposed digit in the source, not a different method.
\end{lrtrapbox}

% =====================================================================
\lo{LTR-12 c}{Discuss factors affecting the choice of non-deposit funding sources.}

Which non-deposit sources management uses to cover a projected available funds gap depends largely on \textbf{five factors}.

\begin{center}
\small
\begin{tabularx}{\textwidth}{@{}p{3.6cm} X@{}}
\toprule
\rowcolor{FRMSoft}\textbf{Factor} & \textbf{What it means} \\
\midrule
\textbf{1. Cost} & Interest rates and other fees associated with the source. Generally \textbf{federal funds are cheapest}, followed by CDs and Eurocurrency deposits; \textbf{commercial paper tends to be more expensive}. \\
\addlinespace
\textbf{2. Risk} & The volatility and dependability of the source. \textbf{Fed funds are most reliable}; CDs and commercial paper are riskier because they depend on market conditions. \\
\addlinespace
\textbf{3. Maturity} & How long the funds are needed. \textbf{Fed funds} for short-term needs; \textbf{long-term debt} for long-term needs. \\
\addlinespace
\textbf{4. Size of the institution} & Smaller banks may not have access to all funding sources, because of minimum investment amounts. \\
\addlinespace
\textbf{5. Regulations} & May limit the amount or type of non-deposit funding a bank can use. \\
\bottomrule
\end{tabularx}
\end{center}

\begin{lrkeybox}[The cost ordering, stated precisely]
Among the cheapest short-term borrowed funds is usually the prevailing \textbf{effective interest rate on federal funds} loaned overnight. Interest rates on \textbf{domestic CDs and Eurocurrency deposits} are in most cases \emph{slightly higher} than the fed funds rate. \textbf{Commercial paper} normally may be issued at rates \emph{slightly above} the fed funds and CD rates, depending on maturity and time of issue. The \textbf{discount rate} on loans from the Federal Reserve banks --- the primary credit rate --- is generally \textbf{among the highest} short-term borrowing rates, because Federal Reserve credit is priced above the central bank's target for the fed funds rate.
\end{lrkeybox}

\begin{lrtrapbox}[The ranking question is a shortcut question]
A stem giving five funding sources and asking for the cheapest is answered by the ordering alone: \textbf{fed funds $<$ domestic CDs $\approx$ Eurocurrency $<$ commercial paper $<$ primary credit}. No arithmetic required, and \S5B.4 grades these on whether the shortcut was found. The counterintuitive end is the right-hand one: \textbf{central bank credit is expensive, by design}.
\end{lrtrapbox}

% =====================================================================
\lo{LTR-12 d}{Calculate overall cost of funds using both the historical average cost approach and the pooled-funds approach.}

\gapbadge
\gapnote{This objective has a heading and no content whatsoever in the source notes. Everything below is written from the GARP curriculum (Rose and Hudgins, Chapter 13).}

\begin{lrgapbox}
The objective is purely computational and names two methods. What is needed is each method's question, its formula chain, and the difference in what each one is for.
\end{lrgapbox}

\begin{lrkeybox}[The two methods answer different questions]
\begin{itemize}
\item The \textbf{historical average cost approach} looks at the \textbf{past}. It asks: what funds has the financial firm raised to date, and what did they cost?
\item The \textbf{pooled-funds approach} looks at the \textbf{future}. It asks: what minimum rate of return must be earned on any \emph{future} loans and investments, just to cover the cost of all \emph{new} funds raised?
\end{itemize}
Backward-looking versus forward-looking is the discriminator, and it is the one GARP tests.
\end{lrkeybox}

\subsection{The historical average cost approach}

\begin{center}
\small
\begin{tabular}{@{}p{5.6cm} p{3.1cm} p{2.7cm} p{2.9cm}@{}}
\toprule
\rowcolor{FRMSoft}\textbf{Source of funds drawn upon} & \raggedright\textbf{Avg.\ amount raised (\$m)} & \raggedright\textbf{Avg.\ interest rate} & \textbf{Total interest paid (\$m)} \\
\midrule
Non-interest-bearing demand deposits & 100 & 0\% & 0 \\
Interest-bearing transaction deposits & 200 & 7\% & 14 \\
Savings accounts & 100 & 5\% & 5 \\
Time deposits & 500 & 8\% & 40 \\
Money market borrowings & 100 & 6\% & 6 \\
\midrule
\textbf{Total} & \textbf{1,000} & & \textbf{65} \\
\bottomrule
\end{tabular}
\end{center}

\begin{lrexbox}[Working the historical average cost, in three steps]
Earning assets are \$750 million, other operating costs are \$10 million, stockholders' investment is \$100 million, the required after-tax return to stockholders is 12\%, and the tax rate is 35\%.
\tcblower
\textbf{Step 1 --- weighted average interest expense.}
\[ \frac{\text{All interest paid}}{\text{Total funds raised}} = \frac{65}{1{,}000} = \mathbf{6.5\%} \]
This is the cost of the money alone. It ignores everything it takes to run the bank.

\textbf{Step 2 --- break-even cost rate on funds invested in earning assets.}
\[ \frac{\text{Interest} + \text{Other operating costs}}{\text{All earning assets}} = \frac{65 + 10}{750} = \mathbf{10\%} \]
It is called \term{break-even} because the institution must earn at least this rate on its earning assets --- primarily loans and securities --- just to meet the total operating costs of raising the funds. Note the \textbf{denominator changes} between step 1 and step 2: total funds raised (\$1,000m) becomes earning assets (\$750m).

\textbf{Step 3 --- weighted average overall cost of capital.}
\[ \text{Break-even cost} \;+\; \frac{\text{After-tax cost of stockholders' investment}}{(1 - \text{Tax rate})} \times \frac{\text{Stockholders' investment}}{\text{Earning assets}} \]
\[ = 10\% \;+\; \frac{12\%}{1 - 0.35} \times \frac{100}{750} \;=\; 10\% + 18.46\% \times 0.1333 \;=\; 10\% + 2.46\% \;=\; \mathbf{12.46\%} \]
\textbf{What this means.} 12.46\% is the lowest rate of return, over all fund-raising costs, that the institution can afford to earn on its assets given \$100 million of equity. The GARP text rounds the middle term to 2.5\% and states the answer as \textbf{12.5\%}; both are the same calculation.
\end{lrexbox}

\begin{lrtrapbox}[Three traps inside the historical average cost chain]
\begin{itemize}
\item \textbf{The denominator switches.} Step 1 divides by \textbf{total funds raised}; steps 2 and 3 divide by \textbf{earning assets}. Using \$1,000m throughout gives $75/1000 = 7.5\%$ at step 2 --- a clean, plausible, wrong number.
\item \textbf{The equity cost is grossed up.} The 12\% is an \textbf{after-tax} required return, so it is divided by $(1 - 0.35)$ to put it on a before-tax basis. Using 12\% directly gives $10\% + 1.6\% = 11.6\%$.
\item \textbf{Intermediate results.} 6.5\% and 10\% are both correct answers to \emph{different} questions in the same chain. \S5B.3's favourite trap is offering the legitimate intermediate quantity, and this calculation supplies two of them.
\end{itemize}
\end{lrtrapbox}

\subsection{The pooled-funds approach}

\begin{center}
\small
\begin{tabular}{@{}p{4.3cm} p{2.1cm} p{2.5cm} p{2.6cm} p{2.5cm}@{}}
\toprule
\rowcolor{FRMSoft}\textbf{New deposits and non-deposit borrowings} & \raggedright\textbf{Amount raised (\$m)} & \raggedright\textbf{Portion placeable in earning assets} & \raggedright\textbf{Dollars into earning assets (\$m)} & \textbf{Operating expenses (\$m)} \\
\midrule
Interest-bearing transaction deposits & 100 & 50\% & 50 & 8 \\
Time deposits & 100 & 60\% & 60 & 9 \\
New stockholders' investment & 100 & 90\% & 90 & 13 \\
\midrule
\textbf{Total} & \textbf{300} & & \textbf{200} & \textbf{30} \\
\bottomrule
\end{tabular}
\end{center}
\figcap{The fourth column is the whole reason the two answers below differ. Only \$200m of the \$300m raised can actually be put to work; the rest is absorbed by reserves and non-earning requirements.}

\begin{lrexbox}[Working the pooled-funds approach, in two steps]
\textbf{Step 1 --- pooled deposit and non-deposit funds expense.}
\[ \frac{\text{All expected operating expenses}}{\text{All new funds expected}} = \frac{30}{300} = \mathbf{10\%} \]
This is what the new money costs, per dollar raised.

\textbf{Step 2 --- hurdle rate of return over all earning assets.}
\[ \frac{\text{All expected operating costs}}{\text{Dollars available to place in earning assets}} = \frac{30}{200} = \mathbf{15\%} \]
\tcblower
\textbf{Why the answer rises from 10\% to 15\%.} Only \textbf{two-thirds} of the expected new funds --- \$200 million out of \$300 million raised --- will actually be available to acquire earning assets. The same \$30 million of cost has to be recovered from a smaller base, so the required return is $\tfrac{3}{2}$ times larger.

\textbf{What this means.} The firm must earn \textbf{at least 15\% before taxes}, on average, on all the new funds it invests, to fully meet its expected fund-raising costs. The \textbf{15\%} figure is the decision-relevant one; 10\% is the cost, not the hurdle.
\end{lrexbox}

\begin{lrtrapbox}[The numerator is the same in both steps; only the denominator moves]
\$30 million appears on top of both fractions. The entire difference is \$300m versus \$200m. A distractor built from ``the portion \emph{not} placeable'' ($30/100 = 30\%$), or from averaging the three placement percentages ($(50+60+90)/3 \approx 66.7\%$, applied to 10\% to give 15\% by a coincidence of this data), will both be offered --- and the second is dangerous because it reaches the right number by the wrong route. The correct weighting is \textbf{dollars into earning assets}, not an unweighted average of percentages.
\end{lrtrapbox}

\begin{lrtrapbox}[Consolidated trap summary — LTR-12]
\begin{itemize}
\item \textbf{Calculation.} $\text{AFG} = (\text{desired loans} + \text{investments} + \text{expected line drawings}) - (\text{current} + \text{expected inflows})$. Expected drawings on committed lines are a \textbf{use}, and expected inflows count as a \textbf{source}.
\item \textbf{Scope.} Institutions add a small buffer on top of the AFG estimate. It is a planning figure, not a floor.
\item \textbf{Sibling.} \textbf{Historical average cost} looks backward; \textbf{pooled funds} looks forward.
\item \textbf{Calculation.} Historical average: weighted interest expense divides by \textbf{total funds raised}; break-even and overall cost of capital divide by \textbf{earning assets}.
\item \textbf{Calculation.} The after-tax equity cost is grossed up by $(1 - \text{tax rate})$ before it is added.
\item \textbf{Intermediate result.} 6.5\%, 10\% and 12.46\% are all correct answers to different questions in one chain; 10\% and 15\% likewise in the pooled-funds chain.
\item \textbf{Sibling.} EuroCD $=$ dollar CD issued \emph{outside} the US. Yankee CD $=$ issued by a \emph{foreign} bank through its \emph{US} branch.
\item \textbf{Polarity.} Discount window: primary \textbf{lowest}, seasonal intermediate, secondary \textbf{highest} --- but the primary credit rate is nonetheless among the \emph{highest} short-term market rates, because it is priced above the fed funds target.
\item \textbf{Scope.} Long-term non-deposit debt is treated as \textbf{secondary capital}, not as funding.
\end{itemize}
\end{lrtrapbox}
